% =====================================================================
%  config/biblio.tex
%  Bibliographie avec biblatex + biber (compatible Unicode/multilingue).
%  Style : auteur-date, courant en traductologie et sciences humaines.
% =====================================================================

\usepackage[
  backend=biber,          % biber : gère l'UTF-8 et le tri multilingue
  style=authoryear,       % (auteur, année)
  sorting=nyt,            % tri : nom, année, titre
  sortlocale=fr_FR,       % tri respectant les règles françaises
  language=french,        % langue de la bibliographie (libellés « In », etc.)
  autolang=none,          % pas de bascule auto : on gère les écritures non
                          % latines à la main via \ar{}, \zh{} (voir .bib).
                          % Évite les avertissements « Language not supported ».
  maxbibnames=99,         % liste tous les auteurs dans la bibliographie
  maxcitenames=2,         % « et al. » au-delà de 2 auteurs dans le texte
  giveninits=true,        % initiales des prénoms
  uniquename=init,
  doi=true, isbn=false, url=true,
]{biblatex}

% Fichier de références.
\addbibresource{biblio/references.bib}

% ---------------------------------------------------------------------
%  Note : au chargement, biblatex peut émettre quelques avertissements
%  « Language 'arabic'/'chinese'/'japanese' not supported » car il n'existe
%  pas de fichier de localisation (.lbx) pour ces langues déclarées à
%  polyglossia. Ces avertissements sont SANS CONSÉQUENCE : la bibliographie
%  est en français et les rares champs en écriture non latine sont gérés à
%  la main via \ar{} / \zh{} dans le fichier .bib.
%
%  IMPORTANT : ne PAS utiliser \DeclareLanguageMapping{chinese}{french}
%  (etc.) pour masquer ces avertissements : combiné à \setotherlanguage
%  {chinese}, cela casse la localisation française de biblatex (libellés
%  « In »/« byeditor » en anglais et sans espace). On préfère donc laisser
%  ces avertissements bénins.

% ---------------------------------------------------------------------
%  Réglages fins conseillés.
% ---------------------------------------------------------------------
% Guillemets français autour des titres d'articles (via csquotes).
\DeclareFieldFormat[article,inbook,incollection,inproceedings]{title}%
  {\enquote{#1}}

% Espace insécable avant les deux-points, etc. : géré par polyglossia
% (langue française) — rien à faire ici.
